\section{Paired Motion and Body Corpus}\label{sec:dataset}

We construct a derived corpus in which a source clip identity connects a motion
description, 128 generated motion variants, and their matching simulation
assets. A \emph{clip} denotes a source segment, including a mirrored segment
when present; a \emph{variant} denotes that clip on one body configuration.
Figure~\ref{fig:data_pipeline} shows the construction paths.

\subsection{Source motion and body configurations}\label{sec:source_data}

\paragraph{Motion.}
We use the AMASS-derived motion portion of HumanML3D~\cite{amass,humanml3d},
drawn from the 18 source collections listed in
Appendix~\ref{app:preprocessing}. Source SMPL+H sequences are processed at
20\,Hz, with clip extraction, temporal cropping, mirroring, and
treadmill/skating cleanup following the HumanML3D/HUMOS pipeline, so
HumanML3D identities are preserved. The pipeline produced 22,459 candidate
clip identities, counting original and mirrored clips; a text-based filter
removed 1,508 that appear to require external objects or environmental
support, leaving 20,951 identities for flat-ground imitation. The filter is a
task-scope heuristic and does not prove that every retained reference is
support-free or dynamically feasible. HumanML3D descriptions stay associated
with the variants of their source clip as metadata; their semantic fidelity is
not audited (Appendix~\ref{app:captions}), and they supply keyword-based motion
classes only for the exploratory analysis in Section~\ref{sec:physics}.

\paragraph{Bodies.}
A body configuration is $b=(g,\boldsymbol{\beta})$, where $g$ identifies the
female or male SMPL variant and $\boldsymbol{\beta}\in\R^{10}$ contains shape
coefficients. We sample
\begin{equation}
    \beta_{ij}\sim\mathcal{U}(-3,3),\qquad
    i=1,\ldots,64,\quad j=1,\ldots,10,
    \label{eq:shape_sampling}
\end{equation}
and cross the 64 vectors with the two variants, giving 128 body configurations
rather than 128 independent coefficient vectors. Uniform sampling is a
controlled source of variation, not a representative sample of human
populations.

\paragraph{Reference and asset generation.}
For each clip, the pretrained HUMOS model~\cite{humos} receives the processed
pose and root features and the target body parameters, and returns a
shape-conditioned reference (root-relative 6D rotations with translation). In
parallel, SMPL bodies~\cite{smpl} with identical parameters are converted to
MJCF assets with SMPLSim~\cite{smplsim}, using mass estimation and hand and
foot simplifications. Each asset has 24 rigid bodies and 69 actuated joint
coordinates, with its own collision proxies, joint ranges, and estimated masses
and inertias. These are simulation-model estimates, not measurements of human
participants.

\subsection{Reference conversion and refinement}\label{sec:refinement}

Generated sequences are converted to ProtoMotions' Z-up convention and
resampled from 20 to 30\,Hz. Forward kinematics (FK) uses the matching asset
rather than a single neutral skeleton, and the packaged representation stores
global body positions and rotations, local rotations, joint positions,
velocities, and contact labels. Refinement is applied independently to each
variant in five steps; default settings are in
Table~\ref{tab:refinement_parameters}.

\emph{Contacts.} Ankle and toe contacts are validated, or detected from foot
position and velocity when labels are unusable; short gaps are closed and very
short runs removed. These label candidate stance intervals and are not measured
forces. \emph{Rotation filtering.} Joint rotations are filtered with
sign-continuous quaternions, normalization, and a bounded geodesic correction;
contact transitions divide the sequence into segments so that filtering does
not average across a change of stance. \emph{Root correction.} A conservative
least-squares root translation, anchored at each contact's median foot
position, reduces horizontal stance sliding. It is temporally regularized and
capped, so it does not guarantee exact foot locking and does not solve for new
joint angles. \emph{Ground clearance.} With floor height zero and margin
$\epsilon_z$, the lift required at frame $t$ is
\begin{equation}
    \ell_t^{\mathrm{req}}
      = \max\!\left(0,\epsilon_z-\min_{x\in\mathcal{C}_b(t)}x_z\right),
    \label{eq:clearance}
\end{equation}
where $\mathcal{C}_b(t)$ is the collision geometry of body $b$, and a smoothed
lift envelope stays above this requirement. This prevents penetration of the
checked geometry in the reference; it does not establish balanced support,
realistic impacts, or freedom from self-collision. \emph{Reconstruction.} FK
and velocities are recomputed from the corrected root and rotations within
each motion.

The objective is modest: reduce measurable kinematic artifacts while preserving
the generated motion. Section~\ref{sec:refinement_results} reports artifact and
fidelity measures; improved kinematics alone does not demonstrate improved
policy learning or dynamic feasibility.

\subsection{Scale and splits}\label{sec:splits}

Each per-clip file stores the 128 variants with frame counts, sampling rates,
and the clip and asset identities needed to recover body parameters, so the
training system can stream clip files without downloading the complete corpus.
Splits are by clip identity: all variants of a clip, and original/mirrored
partners, share a partition (Table~\ref{tab:dataset_counts}), and variants are
generated combinations rather than independent captures. The split is
versioned with seed 42 and roughly 90/5/5 proportions; the 150 development
clips used for design comparisons and the reference audit are reserved within
the training split (Appendix~\ref{app:provenance}), and the test manifest is
never loaded by the training workflow. The splits test held-out clip identities
on the trained body set. They have not been shown to be disjoint by subject or
capture recording, which limits claims about broader motion generalization.

\begin{table}[t]
\centering
\caption{Corpus and split sizes.}\label{tab:dataset_counts}
\begin{tabular}{lrr}
\toprule
Partition & Clip identities & Motion variants \\
\midrule
Training & 18,855 & 2,413,440 \\
Validation & 1,048 & 134,144 \\
Test & 1,048 & 134,144 \\
\midrule
Total & 20,951 & 2,681,728 \\
\bottomrule
\end{tabular}
\end{table}
